\documentclass[10pt,a4paper]{amsart}
\usepackage[margin=19mm]{geometry}
\usepackage{amsmath,amssymb,mathtools}
\usepackage[scr=rsfs,cal=euler]{mathalfa}
\usepackage{xfrac}
\usepackage[unicode,hidelinks,bookmarks=false]{hyperref}
\newcommand{\ie}{i.e.\ }
\newcommand{\cf}{cf.\ }
\newcommand{\resp}{respectively\ }
\newcommand{\Subsection}{\subsection}
\providecommand{\nobreakdash}{\nobreak\hbox{-}}
\setlength{\emergencystretch}{2em}
\multlinegap0pt
\usepackage{algorithm}\usepackage{algpseudocode}
\usepackage{enumerate}
\newtheorem{theorem}{Theorem}
\newcommand{\flag}{\mathrm{flag}}
\title{Probabilistic Floating-Point Round-Off Analysis via Concentration Inequalities}
\author{Yichen Tao, Hongfei Fu, Jiawei Chen, Jean-Baptiste Jeannin}
\date{Source: arXiv:2605.04232v1 (CC BY 4.0)}
\begin{document}
\maketitle

\noindent\emph{Source and licence.} Probabilistic Floating-Point Round-Off Analysis via Concentration Inequalities. Version arXiv:2605.04232v1 (CC BY 4.0), distributed by arXiv under the Creative Commons Attribution 4.0 International licence, http://creativecommons.org/licenses/by/4.0/, as recorded in the arXiv licence metadata for this exact version and shown on the abstract page https://arxiv.org/abs/2605.04232v1. Full licence notice: \texttt{licenses/arxiv-2605.04232-CC-BY-4.0.txt}.\par\smallskip

\noindent\textit{Editorial scope.} Counted unit: the article's theorem thm:partitioncorrect (source0.tex lines 1099--1102). The article's complete original proof of it is delivered with it (source0.tex lines 1104--1125, one \texttt{proof} environment of 22 lines). No mathematics is written, repaired or re-derived: the statement and the proof are the article's own lines, in the article's own notation, and no retained line is substituted or re-flowed.  Retained uncounted as the source-local context the proof needs: lines 821--830; lines 1043--1069.  Nothing was omitted from any retained span, so the package carries no omission record.  No retained line contains a citation, so the wrapper carries no bibliography and no bibliography entry.  No retained line contains a figure environment, an image-inclusion command or drawing code, so no image is delivered and the package declares no asset dependency.  Editorial formatting only, in addition: an AMS \texttt{amsart} preamble that restates the article's own declarations and macros used by the retained lines, namely the environments \texttt{theorem} on separate counters, together with the standard packages those lines need.  The retained environments print in this document as Theorem 1 in order of appearance, whereas the article prints the same items with the numbering of its own full document.  The complete native source of the article as distributed in the arXiv e-print for this exact version is bundled unchanged as \texttt{sources/199841/source0.tex}.
\begin{theorem}
Let $u$ be a fixed threshold.
Let $p(\mathbf x)$ be a function satisfying $F(\mathbf x, \mathbf e) \le p(\mathbf x) \cdot \epsilon$ for all possible values of $\mathbf x$ and $\mathbf e$.
Let $K_u(\mathbf x)\triangleq p(\mathbf x)\cdot \epsilon - u$, and $\mu \triangleq \mathbb E[K_u(\mathbf x)]$. 
Suppose that $\mu < 0$.
For an even analysis order $n$, if we define
$\flag_u \triangleq {\mathbb E[(K_u(\mathbf x) - \mu)^n]}/{\mu^n}$,
then the first-order violation probability admits the following upper bound: $\mathbb P[F(\mathbf x, \mathbf e) \ge u] \le \flag_u$.
\label{thm:cmb-flag}
\end{theorem}

\begin{algorithm}[ht]
\caption{Computing $\flag_u$ with range partition}
\label{alg:partition}
\begin{algorithmic}[1]
\Require Input variables $x_1, \dots, x_m$ and their ranges $I_1, \dots, I_m$, threshold $u$, number of partitions $b$, expression of $K_u(\mathbf x)$ {(either in polynomial or fractional case)}, analysis order $n$
\Ensure An upper bound for first-order violation probability $\flag_u$
\For {$i \gets 1 \text{ to }m$}
    \State Partition range $I_i$ into sub-ranges $I_{i, 1},\dots,I_{i, b}$
    \State Compute weights: $w_{i,j} \gets \mathbb P[x_i\in I_{i,j}]$ for $j = 1$ to $b$
\EndFor
\State $\flag_u \gets 0$
\For {each multi-index $\mathbf i\in \{1, \ldots, b\}^m$}
    \State Compute sub-region weight $W_{\mathbf i} \gets \prod_{j=1}^mw_{j, i_j}$ 
    \State $\mu \gets \mathbb E[K_u(\mathbf x)\mid \mathbf x\in B_{\mathbf i}]$ 
        % \Comment Expectation based on normalized distribution of $\mathbf x$.
    \If {$\mu < 0$}
        \State $\flag_{u, B_{\mathbf i}} \gets \mathbb E[(K_u(\mathbf x) - \mu)^n \mid \mathbf x \in B_{\mathbf i}]/\mu^n$
        % \Comment Subscript of $B$ omitted for brevity.
    \Else 
        \State $\flag_{u, B_{\mathbf i}} \gets 1$
    \EndIf 
    \State Set $\flag_{u, B_{\mathbf i}}$ to $1$ if $\flag_{u, B_{\mathbf i}} > 1$
    \State $\flag_u \gets \flag_u + W_{\mathbf i} \cdot \flag_{u, B_{\mathbf i}}$
\EndFor
\State \textbf{Return} $\flag_u$
\end{algorithmic}
\end{algorithm} 

\begin{theorem}
Let $\flag_u$ be the output of Algorithm~\ref{alg:partition}, then $\mathbb P[K_u(\mathbf x) \ge 0] \le \flag _u$.
\label{thm:partitioncorrect}
\end{theorem}

\begin{proof}
We first establish the correctness of each iteration of the for-loop (line 6 -- 16) by showing that the $\flag_{u, B_{\mathbf i}}$ generated in line 14 is a sound local upper bound 
$\mathbb P[K_u \ge 0 \mid \mathbf x \in B_{\mathbf i}] \le \flag_{u, B_{\mathbf i}}.$
If $\mu > 0$  or the value computed in line 10 exceeds $1$, the algorithm sets $\flag_{u,B_{\mathbf i}}$ to $1$, and the desired inequality holds trivially. 
Otherwise, the validity of the bound follows by an argument analogous to that in Theorem~\ref{thm:cmb-flag}.
Specifically, under the normalized distribution of $\mathbf x$ conditioned on $\mathbf x\in B_{\mathbf i}$,
$$\begin{aligned}
\mathbb P[K_u(\mathbf x) \ge 0\mid \mathbf x \in B_{\mathbf i}] &\le \mathbb P[|K_u(\mathbf x) - \mu| \ge |\mu|\mid \mathbf x \in B_{\mathbf i}] =\mathbb P[(K_u(\mathbf x) - \mu)^n \ge \mu^n\mid \mathbf x \in B_{\mathbf i}] \\
&\le {\mathbb E[(K_u(\mathbf x) - \mu)^n\mid \mathbf x \in B_{\mathbf i}]}/{\mu^n} = \flag_{u, B_{\mathbf i}}.
\end{aligned}$$
This confirms that the local bound $\flag_{u, B_{\mathbf i}}$ indeed bounds the probability. 

Note that the returned $\flag_u$ is equal to the weighted sum of all local bounds, i.e.,
$\flag_u = \sum_{\mathbf i \in (1, \dots, b)^m}\flag_{u, B_{\mathbf i}}W_{\mathbf i}.$
Then we extend the local guarantee to a global one.
Since $\{B_{\mathbf i}\}_{\mathbf i\in (1, \dots, b)^m}$ is a set of disjoint sub-regions covering the domain of $\mathbf x$,
by the law of total probability,
$$\begin{aligned}
\mathbb P[K_u(\mathbf x) \ge 0] &= \textstyle\sum_{\mathbf i \in (1, \dots, b)^m} \mathbb P[K_u(\mathbf x) \ge 0 \mid \mathbf x \in B_{\mathbf i}] \cdot \mathbb P[\mathbf x \in B_{\mathbf i}] 
\le \textstyle\sum_{\mathbf i \in (1, \dots, b)^m} \mathbb  \flag_{u, B_{\mathbf i}}\cdot W_{\mathbf i} = \flag_u.
\end{aligned}$$
\end{proof}

\end{document}
